\section*{Instructions \& Guidelines}
\addcontentsline{toc}{section}{Instructions \& Guidelines}

\begin{noticebox}{\faClipboardList\ Exam Instructions}
  This comprehensive exam is designed to test your depth of knowledge in Distributed Systems, Scalable Machine Learning Pipelines, and Privacy-Preserving Cryptographic Systems. Please adhere to the following rules:
  \begin{itemize}
    \item \textbf{Duration:} This is a 24-hour take-home examination. Your booklet must be uploaded to the portal by the deadline.
    \item \textbf{Format:} All answers must be detailed, academically rigorous, and properly cited.
    \item \textbf{Resources:} You may consult academic publications, textbooks, and documentation. You may not collaborate with any human or utilize generative AI systems for composing your answers.
  \end{itemize}
\end{noticebox}

\vspace{1.0cm}

\subsection*{Honour Code \& Pledge}

By signing below, I certify that I have read the rules governing the Doctoral Candidacy Examination and that the work presented in this booklet is entirely my own. I have not collaborated with any third party, and all citations reference legitimate academic sources.

\vspace{1.5cm}

\noindent
\begin{minipage}[t]{0.45\textwidth}
  \rule{\textwidth}{0.5pt}\\
  \textbf{Sarah Chen}\\
  Doctoral Candidate
\end{minipage}
\hfill
\begin{minipage}[t]{0.45\textwidth}
  \rule{\textwidth}{0.5pt}\\
  \textbf{August 3, 2026}\\
  Date
\end{minipage}

\vspace{1.5cm}

\subsection*{Reference Standards}

All citations in this booklet use the standard \textbf{BibTeX} author-year or numerical scheme and refer to peer-reviewed publications. The corresponding bibliography is located at the end of the booklet. Specific mathematical notations follow standard conventions:
\begin{enumerate}
  \item Matrices are denoted in bold uppercase letters (e.g., $\mathbf{X}$).
  \item State machine states are enclosed in circles in diagrams.
  \item Network messages are represented as solid directed arrows.
\end{enumerate}
